\begingroup
\begin{flushleft}
\bfseries
\fontfamily{phv}
{\Large
Jersey Recognition and Team Affiliation on Soccer Game State Reconstruction}

\mdseries
\end{flushleft}
\normalsize
\let\clearpage\relax
\chapter{ABSTRACT}
\endgroup

\noindent The SoccerNet Game State Reconstruction (GSR) challenge consists of reconstructing, on a
two-dimensional map of the pitch and from single-camera video, the position, team, role and jersey
number of every player. In the reference pipeline, two components stand out for their poor
performance, namely jersey number recognition and the assignment of team and pitch side. This
dissertation intervenes in four of the nine modules of the pipeline, namely it replaces the jersey number
reader, rewrites the team and pitch-side modules and extends re-identification with a colour
descriptor, leaving detection, tracking, calibration and role classification unchanged, so that
any gains can be attributed unambiguously to the changes introduced.

Jersey number recognition was rebuilt around a multimodal vision-and-language model, quantised to
fit within the available graphics memory and applied to back-region crops obtained through pose
estimation, with the successive readings of each tracklet combined by voting. For team assignment,
the appearance-based clustering now audits its own quality, rejecting any partition that leaves
almost every player in a single group, and falls back on jersey colour when
re-identification does not separate the teams, and the pitch side is decided by a cascade of
criteria ending with the goalkeeper's position.

The evaluation covered the $159$ annotated videos of SoccerNet-GSR, comparing the proposed system
with the reference solution run under the same conditions and on the same tracklets. Jersey number
accuracy rose from $15.3\%$ to $37.8\%$, improving on $156$ of the $159$ videos;
team and pitch-side accuracy rose from $83.5\%$ to $92.8\%$, and the number of videos in
which the team module fails, assessed by visual inspection, fell from $28$ to
$7$. GS-HOTA, the official challenge metric, rose from $31.30\%$ to $43.76\%$
on the test set, with gains of the same order on validation and on the training set, which played
no part in development. Decomposing the metric by attribute assigns the gain entirely to the two
attributes that were changed. The comparison also uncovered a calibration defect in the distributed
configuration of the reference solution, which was corrected before comparing.

Controlled experiments show that the remaining margin in jersey number recognition is bounded by the
information present in the images. For team assignment, visual inspection characterised the
remaining failures, most of them due to the unchanged role classifier. Five videos are further shown
to have their team sides inverted in the reference annotation, which penalises systems that get them
right.

\hspace{3cm}\\
\indent
\textbf{Keywords: } Game State Reconstruction, Jersey Number Recognition, Team Affiliation,
Multimodal Models, SoccerNet
